\chapter{Conclusion and Future Work}
\label{ch:conclusion}








%TODO THIS WILL BE MOVED
\section{Change-detection-based transition analysis}
\label{sec:transition_cusum}

The transition analysis presented in \Cref{sec:transition_baseline} uses a
stable-run criterion ($n_{\text{stable}}$ consecutive windows of the new class)
on the per-window classifier output. While this gives a usable detection time,
it relies on the ground-truth target label to define what to look for, and
discards the soft information in the classifier's posterior. To address both
limitations we additionally implement an online change-detection algorithm in
the spirit of \cite[\S6.4]{Blanke2003}. The detector consumes only the
classifier's posterior probability stream; it decides by itself \emph{whether}
a change has occurred and \emph{to which class}, using a principled threshold
tied to a target false-alarm rate.

\subsection{Problem formulation}
\label{sec:cusum_problem}

Following Problem~6.6 of \cite{Blanke2003}, let
$\mathbf{z}(k)\in\Delta^{C-1}$ be the vector of class posteriors emitted by the
per-window classifier at time index $k$, where $C$ is the number of terrain
classes. Let $c_0$ denote the currently active class (the class declared after
the last alarm, or initialised from the first windows of the run). At each
time $k$ we test, for every candidate post-change class $c\neq c_0$, the
hypotheses
\begin{align}
  \mathcal{H}_0 &:\; \mathbf{z}(i)\sim p_{c_0} \quad \text{for } 1\le i\le k, \\
  \mathcal{H}_1^{(c)} &:\;
    \mathbf{z}(i)\sim p_{c_0} \text{ for } i<k_0,\;
    \mathbf{z}(i)\sim p_{c}   \text{ for } i\ge k_0,
\end{align}
where the change time $k_0$ is unknown.

\subsection{CUSUM statistic}
\label{sec:cusum_statistic}

For each candidate class $c$ we accumulate the log-likelihood ratio of the
hypothesis $\mathcal{H}_1^{(c)}$ versus $\mathcal{H}_0$ (cf.\ Eqs.~6.62 and
6.91 of the source):
\begin{equation}
  s_c(k) \;=\; \ln\frac{p_c\!\left(\mathbf{z}(k)\right)}{p_{c_0}\!\left(\mathbf{z}(k)\right)},
  \qquad
  g_c(k) \;=\; \max\!\bigl(0,\; g_c(k-1) + s_c(k)\bigr).
  \label{eq:cusum_recursion}
\end{equation}
Following the recursive CUSUM in \S6.4.2 of \cite{Blanke2003}, the decision
function $g_c(k)$ exhibits negative drift before the change and positive drift
afterwards (cf.\ Figs.~6.7--6.8 in the source). An alarm is raised at the
smallest $k$ for which any of the $C-1$ statistics crosses the threshold $h$:
\begin{equation}
  k_a \;=\; \min\!\Bigl\{ k \,:\, \max_{c\neq c_0} g_c(k) > h \Bigr\},
  \qquad
  \hat{c} \;=\; \arg\max_{c\neq c_0} g_c(k_a).
  \label{eq:cusum_alarm}
\end{equation}
The estimated change occurrence time is recovered as
\begin{equation}
  \hat{k}_0 \;=\; \arg\min_{1\le j\le k_a} S_{\hat{c}}(j),
  \qquad
  S_c(k) \;=\; \sum_{i=1}^{k} s_c(i),
  \label{eq:cusum_kzero}
\end{equation}
i.e.\ the time at which the cumulative sum $S_{\hat{c}}$ last reached its
minimum before the alarm (Eq.~6.63 of the source). After every alarm the
active class is updated to $c_0\leftarrow\hat{c}$ and all statistics are reset
to zero (the reinitialisation step of Algorithm~6.3 in \cite{Blanke2003}).

\subsection{Threshold selection}
\label{sec:cusum_threshold}

The threshold $h$ controls the trade-off between detection delay and
false-alarm rate. Following Wald's approximation reproduced in \S6.4
(cf.\ Remark~6.4 and Eq.~6.74 of the source), the average run length under
$\mathcal{H}_0$ --- i.e.\ the mean time between false alarms (MTBFA) --- grows
approximately exponentially in $h$, while the mean detection delay grows
approximately linearly. We set $h$ on stationary segments of the training runs
to target a chosen MTBFA, then use the \emph{same} $h$ at evaluation time. As
a sanity check we report results for several values of $h$, producing a
delay-vs-false-alarm operating curve that is directly analogous to the
$n_{\text{stable}}$ sweep in \Cref{sec:transition_baseline}.

\subsection{Evaluation against ground-truth transitions}
\label{sec:cusum_evaluation}

The detector emits an unordered list of alarms
$\bigl\{\bigl(k_a^{(j)},\,\hat{c}^{(j)}\bigr)\bigr\}$. To produce metrics
comparable to \Cref{sec:transition_baseline}, each ground-truth transition
$(t_k,\ell_k)$ is paired with the first alarm satisfying
\begin{equation}
  t_k - \delta \;\le\; t_{k_a} \;<\; t_{k+1},
\end{equation}
with grace $\delta = 1\,\text{s}$ to allow detections triggered on the last
window before the change. The pairing yields one of four outcomes per
transition:
\begin{itemize}
  \item \textbf{detected} --- an alarm exists in the window and
        $\hat{c} = \ell_k$;
  \item \textbf{wrong-class} --- an alarm exists but $\hat{c}\neq\ell_k$;
  \item \textbf{missed} --- no alarm exists in the window;
  \item \textbf{boundary} --- insufficient windows remain after $t_k$ to
        observe a stable detection.
\end{itemize}
Detection delay is computed as $t_{k_a} - t_k$ exactly as in
\Cref{sec:transition_baseline}. Alarms that fall outside any expected window
are accumulated into a \emph{false-alarm rate} (alarms per minute of
stationary signal), which the original $n_{\text{stable}}$ evaluation could
not measure.

\subsection{Vector-CUSUM variant on IMU features}
\label{sec:cusum_vector}

An alternative formulation independent of the classifier follows \S6.4.3 of
\cite{Blanke2003} directly. Treating the standardised IMU feature vector
$\mathbf{z}(k)\in\mathbb{R}^{n_z}$ as Gaussian under each class hypothesis,
$\mathbf{z}\sim\mathcal{N}(\boldsymbol{\mu}_c,\mathbf{Q}_c)$, the per-class
log-likelihood becomes
\begin{equation}
\begin{split}
  s_c(\mathbf{z}(k)) \;=\;
    &-\tfrac{1}{2}\bigl(\mathbf{z}(k)-\boldsymbol{\mu}_c\bigr)^{\!\top}
      \mathbf{Q}_c^{-1}\bigl(\mathbf{z}(k)-\boldsymbol{\mu}_c\bigr) \\
    &+\tfrac{1}{2}\bigl(\mathbf{z}(k)-\boldsymbol{\mu}_{c_0}\bigr)^{\!\top}
      \mathbf{Q}_{c_0}^{-1}\bigl(\mathbf{z}(k)-\boldsymbol{\mu}_{c_0}\bigr),
\end{split}
\label{eq:cusum_vector}
\end{equation}
with the same recursion
$g_c(k)=\max\!\bigl(0,\,g_c(k-1)+s_c(\mathbf{z}(k))\bigr)$ and threshold $h$.
The class statistics $(\boldsymbol{\mu}_c,\mathbf{Q}_c)$ are estimated from
the training folds. This is Algorithm~6.6 of the source applied per class,
and is reported as a baseline detector that does not depend on classifier
accuracy.